%% example/text/frontmatter/05_Nomenclature.tex
%% Copyright (c) 2026 Marvin Meck
%
% This work may be distributed and/or modified under the
% conditions of the LaTeX Project Public License, either version 1.3
% of this license or (at your option) any later version.
% The latest version of this license is in
%   https://www.latex-project.org/lppl.txt
% and version 1.3c or later is part of all distributions of LaTeX
% version 2008 or later.
%
% This work has the LPPL maintenance status `author-maintained'.
%
% The Current Maintainer of this work is Marvin Meck.
\begingroup
\renewcommand{\glossarysection}[2][]{}

\chapter*{Nomenclature}

The following summarizes the notation adopted throughout this work and lists the symbols and abbreviations used. 
Standard conventions are followed where possible, and any deviations are introduced where they occur. 

\section*{General notation}

\blindtext[1]

\section*{List of symbols}

The table below summarizes the symbols used in mathematical and physical expressions throughout this work. 
The first column lists the symbols, the second provides a description of their meanings, and the third indicates their dimensions. 
Dimensions are expressed as combinations of the fundamental quantities: length (\(\mathrm{L}\)), mass (\(\mathrm{M}\)), time (\(\mathrm{T}\)), temperature (\(\Theta\)), amount of substance (\(\mathrm{N}\)), and electric current (\(\mathrm{I}\)).

\printunsrtglossary[
  type=symbols,
  style=fstlong3col,
  numberedsection=false,
  nonumberlist=true,
]

\section*{Sub- and superscripts}

\begin{multicols}{2}
  \printunsrtglossary[
    type=scripts,
    style={fstmcol},
    numberedsection=false,
    nonumberlist=true
]
\end{multicols}

\section*{Abbreviations}

\printunsrtglossary[
  type=abbreviations,
  style={fstlong2col},
  numberedsection=false,
  nonumberlist=true,
]

\endgroup
\cleardoublepage
